%\pagestyle{mystyle}
\chapter{Incremental Entity Linking}
% All the papers on knowledge extraction
% \section{IJCKG and NIL Prediction}
% \section{Knowledge Extraction in Specific Domains: Specific Issues}
% \subsection{Legal Judgments}
% \subsection{Instant Messaging Data in Criminal Investigations}
% \section{Organizing Knowledge in Specific Domains}
% Especially infrastructure to organize and hold data.
% \subsection{Judgments}
% \subsection{IM Criminal Investigations}
% \section{Adapting to Specific Domains}

This chapter presents the contributions about \emph{incremental \acl{nel}} (see \Cref{sec:incrementalel}), co-authored during my \phd  These have been published and presented in 2022 at the 11th International Joint Conference on Knowledge Graphs (IJCKG 2022)---\textcite{ijckg2022}---and at the Fourth edition of the Semantic Web Challenge on Tabular Data to Knowledge Graph Matching (SemTab 2022)---\textcite{mammotab}. Both the events were co-located with the 21st International Semantic Web Conference (ISWC 2022).

In 2022, as described in \Cref{sec:incrementalel}, the advancements in the NIL prediction task were left behind the one in \ac{nel}. For instance, the BLINK approach~\cite{blink} was not yet adapted for NIL prediction.
For this reason, we worked towards producing datasets and evaluation procedures for aiding the research in this field, focusing on unstructured text~\cite{ijckg2022} (\Cref{sec:ijckg}), but also on structured tabular data~\cite{mammotab} (\Cref{sec:mammotab})---for which a brief backround review is given in \Cref{sec:sti}.

% \todo{add brief intro with contributions}
% \todo{focus on incremental/nil prediction}

% % abstract from ijckg2022
% In this paper, we
% advocate for batch-based incremental entity extraction methods
% that can exploit NEL with a background KB, detect mentions of
% entities that are not present in the KB yet (NIL mentions), and up-
% date the KB with the novel entities. Based on this assumption, we
% present a methodology to evaluate NEL-based incremental entity
% extraction, which can be applied to a “static” dataset for evaluating
% NEL into a dataset for evaluating incremental entity extraction. We
% apply this methodology to an existing benchmark for evaluating
% NEL algorithms, and evaluate an incremental extraction pipeline
% that orchestrates different strong state-of-the-art and baseline al-
% gorithms for the tasks involved in the extraction process, namely,
% NEL, NIL prediction, and NIL clustering. In presenting our experi-
% ments, we demonstrate the increased difficulty of the information
% extraction task in incremental settings and discuss the strengths of
% the available solutions as well as open challenges.

% % abstract from mammotab
% In this paper, we present MammoTab, a dataset composed of 1M Wikipedia tables extracted from over
% 20M Wikipedia pages and annotated through Wikidata. The lack of this kind of dataset in the state-
% of-the-art makes MammoTab a good resource for testing and training Semantic Table Interpretation
% approaches. The dataset has been designed to cover several key challenges, such as disambiguation,
% homonym and NIL-mentions. The dataset has been evaluated using MTab, one of the best approaches of
% the SemTab challenge.

In particular, in \textcite{ijckg2022}, we focus on batch-based incremental entity extraction, proposing a procedure to adapt existing \ac{nel} datasets to this setting. We apply this method to the WNUM \ac{el} dataset~\cite{Onoe_Durrett_2020} producing its incremental version, \emph{Incremental WNUM}, which is available from GitHub~\footnote{https://github.com/rpo19/Incremental-Entity-Extraction}. This approach leverages \ac{el} with a background \ac{kb}, while also detecting and clustering NIL mentions, and incrementally updating the \ac{kb} with novel entities. We apply this methodology to an existing benchmark, propose a baseline pipeline, and demonstrate the increased difficulty of \acl{ee} in incremental settings.

In \textcite{mammotab}, we introduce \emph{MammoTab}, a large-scale dataset of 1M annotated Wikipedia tables extracted from 20M pages and aligned with Wikidata. This resource was designed to address core challenges in \emph{\acf{sti}}---i.e., the task of identifying entities, classes columns types, and relations between columns in tabular data (see \Cref{sec:sti})---such as entity disambiguation, homonym resolution, and NIL detection.
%Its utility is showcased by evaluating it with one of the leading systems from the SemTab challenge.

Finally, the contributions for incremental \ac{nel} can be summarized as follows:
\begin{enumerate}
    \item A general methodology to adapt any \ac{nel} dataset to incremental \ac{nel}.
    \item A dataset for incremental \ac{el}, suitable for evaluating \ac{nel}, NIL prediction, and NIL clustering---created from an existing \ac{nel} benchmark~\cite{Onoe_Durrett_2020}, applying our methodology.
    \item A baseline pipeline, based on the BLINK bi-encoder~\cite{blink}, which additionally supports NIL prediction and clustering.
    \item Evaluation procedures for incremental \ac{nel}.
    \item An evaluation of the baseline pipeline on the incremental scenario and a discussion of the main challenges.
\end{enumerate}
Similarly the contributions for \ac{sti} are
\begin{itemize}
    \item mammotab with NIL annotations.
    \item evaluation\todo{}
\end{itemize}